% ----- CHAUPAI 1 -----
\subsection*{\deva{प्रथम चौपाई} (First Chaupai)}
\addcontentsline{toc}{subsection}{\deva{प्रथम चौपाई} (First Chaupai) — \deva{जय हनुमान...}}

\begin{minipage}[t]{0.55\textwidth}
\vspace{0pt}

{\large \linespread{1.5}\selectfont
\deva{जय हनुमान ज्ञान गुन सागर।}\\
\deva{जय कपीस तिहुँ लोक उजागर॥}
\par}

\vspace{0.5em}

{\large \linespread{1.5}\selectfont
\telu{జయ హనుమాన జ్ఞాన గున సాగర।}\\
\telu{జయ కపీస తిహుఁ లోక ఉజాగర॥}
\par}

\vspace{0.5em}

{\large \linespread{1.3}\selectfont
\textit{jaya hanumāna jñāna guna sāgara |}\\
\textit{jaya kapīsa tihũ loka ujāgara ||}
\par}
\end{minipage}%
\hfill
\begin{minipage}[t]{0.42\textwidth}
\vspace{0pt}
\begin{tcolorbox}[anchorbox]
\textbf{The Power of Good Communication:} Hanuman is praised here as an "Ocean of Knowledge." When Lord Rama first met him in the Kishkindha forest, Rama was immediately blown away! Not by his physical strength, but by his flawless speech, clear grammar, and humble posture. True intelligence (\textit{Jnana}) is best expressed through respectful, perfectly articulated communication.
\vspace{0.5em}\newline Lord Rama instantly respects Hanuman's flawless communication, proving that meticulous education commands the highest respect.\newline\null\hfill\href{https://www.valmiki.iitk.ac.in/sloka?field_kanda_tid=4&language=dv&field_sarga_value=3}{\textit{Kishkindha Kanda, 3:28-33}}
\end{tcolorbox}
\end{minipage}

\vspace{0.5em}
\subsubsection*{\textbf{Visual Rhythm Map:}}
\begin{table}[h!]
\centering
\renewcommand{\arraystretch}{1.2}
\setlength{\tabcolsep}{0pt}
\begin{tabularx}{\textwidth}{|*{16}{>{\centering\arraybackslash\hsize=1.0\hsize}X|}}
\hline
\multicolumn{2}{|>{\centering\arraybackslash\hsize=2.0\hsize}X|}{\textbf{\laghu{ज}\laghu{य}} \par (2)} &
\multicolumn{5}{>{\centering\arraybackslash\hsize=5.0\hsize}X|}{\textbf{\laghu{ह}\laghu{नु}\guru{मा}\laghu{न}} \par (5)} &
\multicolumn{3}{>{\centering\arraybackslash\hsize=3.0\hsize}X|}{\textbf{\guru{ज्ञा}\laghu{न}} \par (3)} &
\multicolumn{2}{>{\centering\arraybackslash\hsize=2.0\hsize}X|}{\textbf{\laghu{गु}\laghu{न}} \par (2)} &
\multicolumn{4}{>{\centering\arraybackslash\hsize=4.0\hsize}X|}{\textbf{\guru{सा}\laghu{ग}\laghu{र}} \par (4)} \\
\hline
\end{tabularx}
\vspace{-1.5em}
\end{table}
\begin{table}[h!]
\centering
\renewcommand{\arraystretch}{1.2}
\setlength{\tabcolsep}{0pt}
\begin{tabularx}{\textwidth}{|*{16}{>{\centering\arraybackslash\hsize=1.0\hsize}X|}}
\hline
\multicolumn{2}{|>{\centering\arraybackslash\hsize=2.0\hsize}X|}{\textbf{\laghu{ज}\laghu{य}} \par (2)} &
\multicolumn{4}{>{\centering\arraybackslash\hsize=4.0\hsize}X|}{\textbf{\laghu{क}\guru{पी}\laghu{स}} \par (4)} &
\multicolumn{2}{>{\centering\arraybackslash\hsize=2.0\hsize}X|}{\textbf{\laghu{ति}\laghu{हुँ}} \par (2)} &
\multicolumn{3}{>{\centering\arraybackslash\hsize=3.0\hsize}X|}{\textbf{\guru{लो}\laghu{क}} \par (3)} &
\multicolumn{5}{>{\centering\arraybackslash\hsize=5.0\hsize}X|}{\textbf{\laghu{उ}\guru{जा}\laghu{ग}\laghu{र}} \par (5)} \\
\hline
\end{tabularx}
\vspace{-1.5em}
\end{table}

\vspace{1em}

 \textbf{ \deva{सरल-संस्कृतम्}:}\\
 \deva{हे हनुमान्! भवतः जयः भवतु, यः ज्ञानस्य गुणानां च सागरः अस्ति।}\\
 \deva{हे कपीश (वानराणां स्वामिन्)! भवतः जयः भवतु, यः त्रिषु लोकेषु (स्वर्गे, भूमौ, पाताले) प्रकाशमानः (प्रसिद्धः) अस्ति।}\\

 Victory to you, O Hanuman, who is the ocean of knowledge and virtues.\\
 Victory to the Lord of the Monkeys, who illuminates the three worlds.

 \vspace{1em}
\newpage
\textbf{\deva{प्रतिपदार्थः} (Meaning of each word):}
\begin{table}[h!]
\centering
\renewcommand{\arraystretch}{1.2}
\begin{tabularx}{\textwidth}{@{} l l X X @{}}
\toprule
\textbf{Awadhi} & \textbf{Sanskrit} & \textbf{English} & \textbf{Grammar \& Notes} \\
\midrule
\deva{ज्ञान} / \telu{జ్ఞాన}  &  \deva{ज्ञानस्य}  &  Of knowledge  &   \\
\deva{गुन} / \telu{గున}  &  \deva{गुणानाम्}  &  Of virtues/qualities  & \\ 
\deva{सागर} / \telu{సాగర}  &  \deva{सागरः}  &  Ocean  &   \\
\deva{कपीस} \gotchaicon / \telu{కపీస}  &  \deva{कपीशः (कपि+ईश)}  &  Lord of Monkeys  & \\ 
\deva{तिहुँ} / \telu{తిహుఁ}  &  \deva{त्रिषु}  &  In the three  & $tri \rightarrow tihu$ \\ 
\deva{लोक} / \telu{లోక}  &  \deva{लोकेषु}  &  Worlds  &   \\
\deva{उजागर} / \telu{ఉజాగర}  &  \deva{उजागरः (प्रसिद्धः)}  &  Illuminator / Brightness  &   \\
\bottomrule
\end{tabularx}
\end{table}

\begin{tcolorbox}[gotchabox]
\begin{itemize}
    \item \textbf{The \deva{श/ष} $\rightarrow$ \deva{स} Shift:} \deva{कपीश} (Kapeesha) becomes \deva{कपीस} (Kapeesa).
    \item \textbf{The \deva{ण} $\rightarrow$ \deva{न} Shift:} Sanskrit \deva{गुण} with a retroflex N becomes  \deva{गुन} with a dental 'n'. 
    \item \textbf{Schwa Holding:} To maintain the 16-matra beat, do not cut the final consonants dead. Hold the "r" in Saaga-r and Ujaaga-r for exactly one beat.
\end{itemize}
\end{tcolorbox}

\newpage
